\section{Witness degree and specialization of orbit closures}
Fix integers $1\le k\le N$.  In the affine space $W$ of degree-$N$
forms in $N^2$ variables, let $y=\det_N$ and let
$z=\ell^{N-k}\operatorname{perm}_k$ with a fixed integral embedding
of its variables.  Put $G=\mathrm{GL}_{N^2}$, the general linear group acting by linear changes of variables.  For a field $K$, let
\begin{equation}
 Y_K=\overline{G_{\overline K}\cdot y_K}\subset W_{\overline K},
 \qquad
 c_{k,N}(K)=\min\{\deg f:f\in I(Y_K),\ f(z_K)\ne0\},
\label{eq:WD1}
\end{equation}
with minimum $\infty$ if the set is empty.  This is degree in the
coefficient coordinates on $W$, not degree of the permanent itself.
The same definition applies to any fixed integral representation and
two integral points.  We write $c_0$ for characteristic zero and $c_p$
for characteristic $p$.

The algebraic closure in \eqref{eq:WD1} is essential: vanishing on the finite set
$G(\mathbb F_p)\cdot y$ would define a different problem.  However,
a minimum-degree witness may be chosen over $K$ itself.  Indeed its
existence is the rank condition below, which is unchanged by field
extension.  In the determinant representation $Y_K$ is a cone:
scalar matrices multiply a degree-$N$ form by an $N$th power, and these
powers exhaust $\overline K^*$ even when $p\mid N$.  Thus its ideal is
homogeneous, and a minimum-degree witness can also be chosen homogeneous.

\begin{lemma}[An exact bounded-degree matrix]
For every fixed $D$, there is an integer matrix $M_D$ and an integer row
$e_D$ such that over every field $K$,
\begin{equation}
 c_{k,N}(K)\le D
 \quad\Longleftrightarrow\quad
 \operatorname{rank}_K\begin{pmatrix}M_D\\e_D\end{pmatrix}
 >\operatorname{rank}_K M_D.
\label{eq:WD2}
\end{equation}
\end{lemma}
\begin{proof}
Use the substitution orbit map $g\mapsto\det_N(gx)$.  It is a
polynomial map from the space of all $N^2\times N^2$ matrices, with
integer coefficients, and its restriction to $G$ has the required
orbit.  Since $G_{\overline K}$ is dense in matrix space over every
field, a polynomial vanishes on the orbit if and only if its pullback
is the zero polynomial in the matrix entries.

Index the columns by monomials in the coefficient coordinates of degree
at most $D$.  Their pullbacks have degree at most $ND$ in the matrix
entries.  Expand them in that finite monomial basis to form $M_D$.
Let $e_D$ evaluate the original monomials at $z$.  A vector $v$ of
coefficients defines a witness precisely if $M_Dv=0$ and $e_Dv\ne0$.
Over any field this is equivalent to $e_D$ not belonging to the row
space of $M_D$, which is \eqref{eq:WD2}.  All coordinates are integral.

For a general regular $\mathrm{GL}_r$ representation, its orbit map
has entries in $\mathbb Z[g_{ij},\det(g)^{-1}]$.  For a fixed $D$,
one common power of $\det(g)$ clears every denominator.  Multiplication
by that power is injective over every field, so the identical argument
gives a finite integer coefficient matrix.
\end{proof}

\begin{samepage}
\begin{theorem}[Specialization of the minimum witness degree]
For every fixed pair $(k,N)$:
\begin{enumerate}
\item If $c_0=C<\infty$, then $c_p=C$ for all but finitely many primes.
\item If $c_0=\infty$, then for every fixed $D$, one has $c_p>D$
for all but finitely many primes.  Thus $c_p\to\infty$ in the
extended sense, allowing $c_p=\infty$.
\end{enumerate}
\end{theorem}
\end{samepage}
\begin{proof}
The rank of a fixed integer matrix equals its rational rank modulo
every prime not dividing one chosen nonzero maximal-rank minor.
All minors of larger size vanish as integers.  Applying this to both
matrices in \eqref{eq:WD2} proves that the bounded-degree existence predicate
agrees with its characteristic-zero value outside a finite, explicitly
certifiable set of primes.

If $C<\infty$, apply the assertion at $D=C$ and $D=C-1$ (or note
directly that nonzero constants cannot vanish on a nonempty orbit).
If $c_0=\infty$, apply it separately to each fixed $D$.  The exceptional
set may depend on $D$; exchanging these quantifiers would be invalid.
\end{proof}

There is a simpler certificate for the lower bound at a given $D$.
When no degree-$D$ witness exists in characteristic zero, solve
$e_D=\lambda M_D$ over $\mathbb Q$.  Clearing denominators gives
$b_De_D=\mu M_D$ with $b_D\ne0$ integral.  At any prime not dividing
$b_D$ this proves the absence of such a witness.  Conversely, clearing
denominators in one characteristic-zero witness gives an integral
polynomial whose pullback is identically zero; it remains a witness at
every prime not dividing its nonzero integer value at $z$.  These are explicit certificates for survival and bounded-degree specialization.

\begin{theorem}[Stronger stabilization for fixed orbit data]
For a fixed integral $\mathrm{GL}_r$ representation and fixed integral
points $y,z$, one has $c_p=c_0$ outside a finite set of primes, with
equality also interpreted in $\mathbb N\cup\{\infty\}$.
\end{theorem}
\begin{proof}
Only the case $c_0=\infty$ remains.  We give an explicit
scheme-image argument using elimination; in particular we do not assume
that taking an orbit closure commutes with every reduction.

Let $w_1,\ldots,w_M$ be the coordinates of the representation and write
the orbit map with coordinates in
$\mathbb Z[g_{ij},u]/(u\det(g)-1)$.  Choose polynomial representatives
$P_i(g,u)$ and form the graph ideal
\[
 J=(u\det(g)-1,\ w_i-P_i(g,u):1\le i\le M)
 \subset\mathbb Z[g_{ij},u,w_1,\ldots,w_M].
\]
Over every field this graph has coordinate ring $K[G]$, a domain;
its elimination ideal $J_K\cap K[w]$ is therefore the prime ideal
of the geometric orbit closure.

Over $\mathbb Q$, compute a finite monic Groebner basis $\mathcal G$
with an elimination order in which all $g_{ij},u$ precede the $w_i$.
Retain the finite polynomial identities expressing its members in
terms of the displayed generators of $J$, the generators in terms of
$\mathcal G$, and the reductions to zero of its finitely many
$S$-polynomials.  Invert one nonzero integer $b$ clearing every
denominator in these identities and in $\mathcal G$.  Over
$R=\mathbb Z[1/b]$ they show $J_R=(\mathcal G)$.

For every $p\nmid b$, all leading coefficients of $\mathcal G$ remain
one, its leading monomials are unchanged, and the same $S$-polynomial
reduction identities hold.  Buchberger's criterion gives a Groebner
basis $\mathcal G_p$ of $J_{\mathbb F_p}$.  The elimination theorem
therefore gives
\begin{equation}
 I(Y_{\overline{\mathbb F}_p})
   =(\mathcal H\bmod p),\qquad
 \mathcal H=\mathcal G\cap R[w],
\label{eq:WD3}
\end{equation}
after scalar extension.  The elimination order and persistence of
leading monomials ensure that no member involving an eliminated
variable suddenly becomes a $w$-only polynomial.  The same set
$\mathcal H$ generates the characteristic-zero orbit ideal.

When $c_0=\infty$, the point $z$ belongs to the characteristic-zero
orbit closure, so $h(z)=0$ in $\mathbb Q$ for every $h\in\mathcal H$.
These equalities already hold in $R$ and hence in every residue field
with $p\nmid b$.  Equation \eqref{eq:WD3} gives $z_p\in Y_p$ and $c_p=\infty$
for all such primes.  The finite-$c_0$ case was proved above.
\end{proof}

\paragraph{A second geometric proof of the closure comparison.}
Let $\mathcal Y$ be the schematic image of $G_{\mathbb Z}\to W_{\mathbb Z}$.
Its coordinate ring is the subring of $\mathbb Z[G]$ generated by the
finitely many coordinate pullbacks from $W$. It is therefore a
finite-type integral ring, torsion-free and hence flat over
$\mathbb Z$.  Its generic fiber is geometrically integral, since it
is the closure of the image of the geometrically integral group
$G_{\mathbb Q}$.  By \citet[Theorem~9.7.7(iv), pp.~79--81]{EGAIV3},
the geometrically integral fiber locus is constructible, hence contains
an open neighborhood of the generic point of $\operatorname{Spec}\mathbb Z$.
By Chevalley's theorem \citep[Theorem~1.8.4, p.~239]{EGAIV1}, the image
of $G$ in $\mathcal Y$ is constructible and contains a nonempty open
subset $U$.  This open subset meets the generic fiber.  Applying
Chevalley to $U\to\operatorname{Spec}\mathbb Z$ shows that $U_p$ is
nonempty after excluding finitely many more primes.  For each remaining
prime, $U_p$ is dense in the geometrically integral fiber $\mathcal Y_p$,
and is contained in the image of $G_p$.  After algebraic closure the
same holds: nonempty finite-type fibers over an algebraically closed
field have rational points.  Therefore the fiberwise geometric orbit
closure is exactly $\mathcal Y_p$, including its reduced scheme structure.
This proves the comparison used in \eqref{eq:WD3} independently of elimination.
No assertion of closure commutation at an excluded prime is made.

\paragraph{An exceptional prime and the model distinction.}
Let $\mathbb G_m$ act on $\mathbb A^1_{\mathbb Z}$ by scalar
multiplication, with $y=2$ and $z=1$.  In characteristic zero the orbit
of $2$ is dense in the affine line, so $c_0=\infty$.  Modulo two its
orbit is $\{0\}$, and $f(w)=w$ separates $z=1$, so $c_2=1$.
For every odd prime the orbit is again dense, so $c_p=\infty$.
The integral scheme-theoretic image is already the whole affine line,
which is flat over $\mathbb Z$; its characteristic-two fiber is not the
fiberwise orbit closure.  Flatness of the chosen closure model alone
does not remove this issue.  The finite localization in \eqref{eq:WD3} does.
Thus a single-prime witness need not lift to a characteristic-zero separation.

\paragraph{Scope as a complexity statement.}
The conclusion is for each fixed size pair $(k,N)$.  Neither the degree
$C$, the exceptional integer $b$, nor the size of the coefficient
matrices is bounded polynomially in growing $(k,N)$ by this argument.
Consequently a characteristic-zero separation premise implies finite
$c_{k,N}$ and almost-everywhere equality for each covered pair, but does
not give a uniform small witness, a uniform exceptional-prime bound,
or a circuit-size bound for a witness.  Homogeneous highest-weight
equations separating at an orbit translate may be used in characteristic
zero when assumed, but their highest-weight form and
semisimplicity need not specialize.  The theorem uses ordinary
separating polynomials and remains valid without those additional
representation-theoretic assertions.
